
% Table 1: Main Results
\begin{table}[t]
\centering
\begin{tabular}{lccc}
\hline
\textbf{Model} & \textbf{SR (\%)} & \textbf{SPL (\%)} & \textbf{Oracle SR (\%)} \\
\hline
VLN-BERT (Paper) & ~55 & ~50 & - \\
\textbf{Ours (LLM+ResNet)} & \textbf{62.0} & \textbf{61.8} & \textbf{69.0} \\
\hline
\end{tabular}
\caption{Main results comparison with VLN-BERT baseline.}
\label{tab:main_results}
\end{table}

% Table 2: Instruction Quality Comparison
\begin{table}[t]
\centering
\begin{tabular}{lcc}
\hline
\textbf{Method} & \textbf{Excellent Rate (\%)} & \textbf{Avg Score (1-5)} \\
\hline
Machine Translation & 20.0 & 4.28 \\
\textbf{LLM Direct (Ours)} & \textbf{76.7} & \textbf{4.61} \\
\hline
\end{tabular}
\caption{Instruction quality comparison between machine translation and LLM direct generation.}
\label{tab:instruction_quality}
\end{table}

% Table 3: Feature Ablation Study
\begin{table}[t]
\centering
\begin{tabular}{lcc}
\hline
\textbf{Feature Type} & \textbf{Best Val Loss} & \textbf{Training Time} \\
\hline
Random Gaussian & 3.37 & 6 min \\
\textbf{ResNet-50 (Ours)} & \textbf{2.63 (-22\%)} & \textbf{1 min (-83\%)} \\
\hline
\end{tabular}
\caption{Ablation study on visual feature representations.}
\label{tab:feature_ablation}
\end{table}

% Table 4: Performance by Instruction Type
\begin{table}[t]
\centering
\begin{tabular}{lcc}
\hline
\textbf{Instruction Type} & \textbf{Samples} & \textbf{SR (\%)} \\
\hline
Turn Left & 104 & 58.7 \\
Turn Right & 96 & 65.6 \\
Go Straight & 109 & 68.8 \\
\hline
\end{tabular}
\caption{Performance breakdown by instruction type.}
\label{tab:instruction_type}
\end{table}

% Table 5: Training Configuration
\begin{table}[t]
\centering
\begin{tabular}{lc}
\hline
\textbf{Hyperparameter} & \textbf{Value} \\
\hline
Batch Size & 16 (effective 32) \\
Learning Rate & 1e-4 (warmup + cosine) \\
Warmup Epochs & 3 \\
Max Epochs & 20 \\
Early Stop Patience & 5 \\
Vocabulary Size & 97 chars \\
Hidden Dimension & 256 \\
Attention Heads & 8 \\
Encoder Layers & 2 \\
\hline
\end{tabular}
\caption{Training hyperparameters.}
\label{tab:hyperparams}
\end{table}
